\section{Worked Example: Automated Funding Program Matching}

We illustrate how the tactics can be applied in a concrete deployment: an automated funding program matching tool designed to support citizen-driven innovation. Public funding landscapes are fragmented across thousands of programs whose eligibility criteria are expressed in natural language; applicants---individuals, small organisations, community initiatives---must navigate this complexity without requiring legal expertise. The design goal was to automate eligibility matching while maintaining oversight at points of legal or epistemic uncertainty. The system separates two concerns: an offline stage that compiles natural-language program criteria into executable eligibility rules, and a query stage that matches a project description submitted by the applicant against those rules at runtime.

In the offline stage, an agent writes executable eligibility rules to the program database, placing the system at Agency~L4. \emph{Tool Provisioning}~(T4) equipped the query stage with a geodata API that resolves postcodes and place names into standardised geographic identifiers. Before the search proceeds, a \emph{Checkpoint}~(T1) presents the applicant with the structured project profile derived from their project description---funding amount, location, project type---for review and correction, ensuring the search correctly reflects the applicant's intent. \emph{Escalation}~(T2) handles programs for which a clear eligibility verdict cannot be reached: rather than forcing a decision, the system presents these programs to the applicant with an explicit indication that eligibility could not be confirmed, leaving the final judgement to them. \emph{Tool Fencing}~(T5) constrains the generated eligibility rules to a sandboxed execution environment (no network access, read-only filesystem, dropped privileges), ensuring that LLM-generated logic cannot affect external systems.

